\chapter{Planificación y evolución}

\section{Planificación inicial heredada del GEP}
El GEP definió una planificación de 545 horas entre el 9 de febrero y la última semana de junio de 2026. La estructura era adecuada y se mantiene como base de seguimiento porque separaba seis bloques bien delimitados: gestión y planificación, análisis y requisitos, diseño de arquitectura, implementación e integración, pruebas y calidad, y documentación final.

\begin{table}[H]
  \centering
  \caption{Planificación inicial del GEP}
  \label{tab:planificacion-inicial-gep}
  \begin{tabularx}{\textwidth}{lXr}
    \toprule
    Código & Bloque & Horas \\
    \midrule
    GP & Gestión y planificación & 70 \\
    AN & Análisis y requisitos & 45 \\
    DS & Diseño de arquitectura & 60 \\
    IM & Implementación e integración & 220 \\
    QA & Pruebas y control de calidad & 50 \\
    DO & Documentación final y defensa & 100 \\
    \midrule
    TOT & Total previsto & 545 \\
    \bottomrule
  \end{tabularx}
\end{table}

\section{Desglose de tareas revisado}
\begin{longtable}{p{0.09\textwidth}p{0.27\textwidth}p{0.13\textwidth}p{0.43\textwidth}}
  \caption{Desglose revisado de tareas respecto al GEP}
  \label{tab:tareas-revisadas}\\
  \toprule
  ID & Tarea & Estado & Revisión durante el seguimiento \\
  \midrule
  \endfirsthead
  \toprule
  ID & Tarea & Estado & Revisión durante el seguimiento \\
  \midrule
  \endhead
  \bottomrule
  \endfoot
  GP1 & Contextualización y alcance & Cerrada & Se concreta mejor el alcance real y se excluyen correos y contraseñas. \\
  GP2 & Planificación temporal & Revisada & Se incorporan las desviaciones reales de voz, observabilidad y reorganización. \\
  GP3 & Gestión económica y sostenibilidad & Revisada & Se distinguen presupuesto de referencia y coste incremental real. \\
  AN1 & Estudio hardware y viabilidad & Avanzada & La GPU AMD y ROCm han exigido más análisis del previsto, especialmente en TTS. \\
  AN2 & Análisis modelos LLM & Avanzada & Ollama queda consolidado como servidor de inferencia local. \\
  AN3 & Diseño seguridad perimetral & Reorientada & Tailscale sustituye a WireGuard manual por razones operativas. \\
  DS1 & Topología de contenedores & Revisada & La arquitectura real se reorganiza por dominios funcionales. \\
  DS2 & Diseño orquestador y RAG & Avanzada & n8n, Nextcloud, Qdrant e ingesta ya forman memoria documental funcional. \\
  DS3 & Diseño sensórico y accesos & Avanzada & Home Assistant se mantiene en Raspberry Pi 5 y se integran dispositivos ya disponibles. \\
  IM1 & Sprint 1 - Infraestructura & Avanzada & Servidor, Compose, almacenamiento, Tailscale y servicios base desplegados. \\
  IM2 & Sprint 2 - Cognitivo & Avanzada & Ollama funcional y herramientas de memoria conectadas. \\
  IM3 & Sprint 3 - Orquestación & Avanzada & Tools de Open WebUI y webhooks de n8n convierten conversación en acciones. \\
  IM4 & Sprint 4 - Sensorial & Avanzada & Domótica, cámaras, HA Voice y presencia ya disponibles. \\
  IM5 & Sprint 5+ - Expansión & Parcial & Fotos, monitorización y autoreparabilidad implementadas; dashboard único y backup final pendientes. \\
  QA1 & Estrés y límites VRAM & Parcial & Límites prácticos identificados; faltan métricas finales normalizadas. \\
  QA2 & Auditoría de seguridad & Parcial & Acceso remoto y secretos revisados; falta cierre formal. \\
  DO1 & Memoria técnica & En curso & Ya existe base extensa, bibliografía real y documentación de herramientas. \\
  DO2 & Defensa & Pendiente & Reservada al cierre formal del TFG. \\
\end{longtable}

\section{Cambios producidos}
\begin{itemize}
  \item \textbf{Acceso remoto}: Tailscale sustituye a WireGuard manual y DuckDNS.
  \item \textbf{Voz}: el esfuerzo en TTS/STT se alarga por compatibilidad AMD y latencia.
  \item \textbf{Infraestructura}: se añade reorganización de Docker y saneamiento del repositorio.
  \item \textbf{Monitorización}: se incorporan Netdata, Telegram y exportación a Home Assistant.
  \item \textbf{Autorreparabilidad}: aparece una línea nueva a partir de incidencias reales de Ollama.
\end{itemize}

\section{Impacto en objetivos y costes}
El análisis de impacto sigue el mismo criterio del GEP: una desviación debe relacionarse con objetivos, esfuerzo y presupuesto. El objetivo general no cambia. La evolución real confirma que los riesgos iniciales estaban bien elegidos, especialmente la incompatibilidad AMD/ROCm y el crecimiento del alcance. Por eso, el mayor sobreesfuerzo no exige redefinir el proyecto, sino consumir parte de la reserva técnica ya prevista.

\begin{table}[H]
  \centering
  \caption{Impacto de los cambios sobre objetivos y costes}
  \label{tab:impacto-objetivos-costes}
  \begin{tabularx}{\textwidth}{lXX}
    \toprule
    Cambio & Impacto en objetivos & Impacto económico \\
    \midrule
    Tailscale en lugar de WireGuard manual & Mantiene el acceso remoto seguro con mejor estabilidad. & Sin licencia adicional y con menos mantenimiento. \\
    Mayor esfuerzo en voz & Conserva voz funcional; la excelencia sonora pasa a límite aceptado. & Consume margen asociado a AMD/ROCm. \\
    Reorganización Docker & Mejora mantenibilidad y reproducibilidad. & Añade horas ahora, reduce riesgo futuro. \\
    Monitorización y autoreparabilidad & Refuerza proactividad y operación real. & Reasigna tiempo desde expansión secundaria hacia calidad. \\
    Dashboard y backup final pendientes & No afectan al núcleo funcional alcanzado. & El NVMe externo queda como ampliación pendiente. \\
    \bottomrule
  \end{tabularx}
\end{table}

El presupuesto inicial del GEP se mantiene como referencia: 10.305,67 euros de costes de personal, 1.488,44 euros de costes genéricos, 1.769,12 euros de contingencia y 1.045,11 euros de imprevistos, hasta un total de 14.608,34 euros. El seguimiento no obliga a elevar esa cifra porque no se han incorporado licencias ni compras obligatorias nuevas; se ha redistribuido esfuerzo dentro del mismo marco.

\section{Gantt revisado}
El diagrama inicial era válido como hoja de ruta, pero demasiado esquemático para un seguimiento. La \cref{fig:gantt-seguimiento} conserva los códigos originales y muestra la convivencia real entre fases a lo largo de veinte semanas.

\begin{figure}[H]
  \centering
  \resizebox{\textwidth}{!}{%
    \begin{ganttchart}[
      hgrid,
      vgrid,
      x unit=0.52cm,
      y unit chart=0.54cm,
      bar height=0.46,
      group height=0.28,
      group peaks height=0.18,
      title height=1,
      bar label font=\scriptsize,
      group label font=\scriptsize\bfseries,
      milestone label font=\scriptsize,
      title label font=\scriptsize\bfseries
    ]{1}{20}
      \gantttitle{Febrero}{3}
      \gantttitle{Marzo}{4}
      \gantttitle{Abril}{5}
      \gantttitle{Mayo}{4}
      \gantttitle{Junio}{4} \\
      \gantttitlelist{1,...,20}{1} \\
      \ganttgroup{GP Gestión y planificación}{1}{20} \\
      \ganttbar{GP1 Contexto y alcance}{1}{3} \\
      \ganttbar{GP2 Planificación temporal}{2}{4} \\
      \ganttbar{GP3 Economía y sostenibilidad}{3}{5} \\
      \ganttgroup{AN Análisis y requisitos}{2}{9} \\
      \ganttbar{AN1 Hardware y viabilidad}{2}{7} \\
      \ganttbar{AN2 Modelos LLM}{3}{8} \\
      \ganttbar{AN3 Seguridad perimetral}{4}{9} \\
      \ganttgroup{DS Diseño de arquitectura}{4}{10} \\
      \ganttbar{DS1 Topología contenedores}{4}{8} \\
      \ganttbar{DS2 Orquestación y RAG}{5}{9} \\
      \ganttbar{DS3 Sensorial y accesos}{6}{10} \\
      \ganttgroup{IM Implementación e integración}{6}{17} \\
      \ganttbar{IM1 Infraestructura}{6}{9} \\
      \ganttbar{IM2 Despliegue cognitivo}{7}{11} \\
      \ganttbar{IM3 Orquestación}{9}{13} \\
      \ganttbar{IM4 Sensorial / PMV}{10}{14} \\
      \ganttbar{IM5 Expansión asistente}{12}{17} \\
      \ganttgroup{QA Pruebas y calidad}{10}{19} \\
      \ganttbar{QA1 Estrés y VRAM}{10}{16} \\
      \ganttbar{QA2 Seguridad y operación}{12}{19} \\
      \ganttgroup{DO Documentación final}{1}{20} \\
      \ganttbar{DO1 Memoria técnica}{1}{20} \\
      \ganttbar{DO2 Defensa}{18}{20} \\
      \ganttmilestone{M1 Inicio}{3} \\
      \ganttmilestone{M2 Infraestructura}{9} \\
      \ganttmilestone{M3 Inteligencia}{11} \\
      \ganttmilestone{M4 PMV}{14} \\
      \ganttmilestone{Seguimiento}{14} \\
      \ganttmilestone{M5 Entrega}{20}
    \end{ganttchart}
  }
  \caption{Diagrama de Gantt revisado del proyecto}
  \label{fig:gantt-seguimiento}
\end{figure}

\section{Punto actual}
El seguimiento se sitúa tras alcanzar el PMV funcional. Infraestructura, inferencia, memoria, domótica y monitorización están avanzadas; voz, backup, métricas finales y documentación permanecen como frentes de cierre. El calendario restante debe concentrarse en consolidación, no en abrir nuevas ramas funcionales.
